\chapter{UML tools} \label{ch:procedure}


\section{Draw.io}




UML (Unified Modeling Language) is a visual language used in software engineering to design and understand systems before development. Common UML diagrams include class diagrams, use case diagrams, sequence diagrams, and activity diagrams. Tools like Draw.io are widely used to create these diagrams.
\subsection{Advantages}
\begin{enumerate}
    \item Completely free
    \item Beginner-friendly
    \item No account required
    \item Works both online and offline
    \item Provides many UML shapes and templates
\end{enumerate}

\subsection{Disadvantages}
\begin{enumerate}
    \item Requires internet for online use
    \item Fewer advanced features compared to professional UML tools
    \item Large diagrams may slow down performance
\end{enumerate}

\section{SmartDraw}

SmartDraw is a professional diagramming tool used for creating a wide range of diagrams, including UML diagrams, flowcharts, organizational charts, and system design models. It is commonly used in both academic and professional environments.

\subsection{Advantages}
\begin{enumerate}
    \item Provides a large collection of pre-built UML templates and standardized symbols for quick and consistent diagram creation
    \item Accessible through web-based and desktop versions
    \item Supports integration with cloud services and team collaboration tools
    \item Offers multiple export formats including PNG, PDF, SVG, and integration with Microsoft Office
\end{enumerate}

\subsection{Disadvantages}
\begin{enumerate}
    \item Not completely free; full functionality requires a paid subscription
    \item Some advanced features may feel complex for beginners
    \item Performance issues may occur with very large or detailed diagrams
\end{enumerate}

\section{Lucidchart}

Lucidchart is a cloud-based diagramming platform that enables users to create a variety of visuals, including UML diagrams, flowcharts, wireframes, and network diagrams. It integrates seamlessly with popular productivity suites like Google Workspace, Microsoft Office, and Confluence.

\subsection{Advantages}
\begin{enumerate}
    \item Easy to use with drag-and-drop interface and templates
    \item Excellent for real-time team collaboration
    \item Cloud-based, accessible from anywhere
    \item Integrates well with tools like Google Drive, Slack, and Jira
    \item Offers a versatile template library with 1000+ templates across categories
    \item Provides unlimited undo/redo with version tracking
\end{enumerate}

\subsection{Disadvantages}
\begin{enumerate}
    \item Requires internet connection; limited offline use
    \item Can slow down with very large or complex diagrams
    \item Advanced features require higher-tier paid plans
    \item Performance may vary across different browsers
    \item Less powerful for highly specialized or detailed technical diagrams
\end{enumerate}

\section{PlantUML}

PlantUML is a text-based UML diagramming tool that allows users to write diagrams using a simple descriptive language. It is widely considered the industry standard for "Diagrams-as-Code."

\subsection{Advantages}
\begin{enumerate}
    \item Text-based files integrate well with version control systems like Git
    \item Typing relationships (e.g., A -> B) is faster than manual drag-and-drop
    \item Automatic layout handling saves time by managing alignment and positioning
    \item Works natively in IDEs such as VS Code and IntelliJ, and documentation tools like Confluence
\end{enumerate}

\subsection{Disadvantages}
\begin{enumerate}
    \item Limited manual control over diagram layout and positioning
    \item Default visuals are basic and not presentation-ready
    \item Requires learning specific syntax, which may be a barrier for non-coders
\end{enumerate}

\section{diagrams.net}

diagrams.net is a commonly used online diagramming tool that supports the creation of UML diagrams and other system design diagrams. It is frequently used by students due to its free availability and ease of use.

\subsection{Advantages}
\begin{enumerate}
    \item Free to use, suitable for students and small project teams
    \item Accessible directly through a web browser without installation
    \item Provides a simple drag-and-drop interface for UML diagram creation
    \item Supports multiple diagram types and offers export options such as PNG, PDF, and SVG
    \item Integrates with cloud storage services like Google Drive for easy sharing and collaboration
\end{enumerate}

\subsection{Disadvantages}
\begin{enumerate}
    \item Does not strictly enforce UML standards; users must ensure diagram correctness
    \item Lacks advanced UML features such as automatic code generation and reverse engineering
    \item Managing large or complex diagrams can be challenging
    \item Collaboration features are basic compared to specialized UML tools
\end{enumerate}
